\chapter{正典奶龙：一个原创 IP 的产业史（2017—2026）}\label{ch:03}

\begin{nlkeypoints}
  \begin{enumerate}[label=(\arabic*),leftmargin=2.4em]
    \item 正典奶龙的成长是一条完整的\keyw{内容产业}轨迹：
          2017 年公司成立、2019 年形象构思、2020 年 1 月 22 日上线、
          剧集化与大屏化（2023—2026）、六周年仪式（2026 年 1 月）、
          以及授权体系的规模化（200+ 品牌、2000+ SKU、57 国）。
    \item 对同一指标，本书坚持\keyw{来源口径并列}：
          “全网粉丝超 3800 万”（A 级）与“全球粉丝突破 3900 万+”（B 级）
          并存呈现，不取均值、不互相覆盖。
    \item 正典奶龙是\keyw{子供向}作品，但其受众外溢到高校人群
          是一个真实且可分析的现象：本章提出“内容内核—平台结构—
          民间再生产”三层溢出解释，并在第九章与奶蛙产业线相对照。
  \end{enumerate}
\end{nlkeypoints}

第二章给对象结构划了界：奶龙学的核心是“正典—变种”这一对。
本章先处理这一对中稳定的一端——正典奶龙。
它的重要性有两个：其一，奶蛙的一切辨识度都以它为参照系，
不了解正典，就无法度量变种的“偏离”；
其二，正典奶龙是一个有专业法务团队、有明确授权体系的产业主体，
它与奶蛙之间的张力（关评、拉黑、维权）构成本书
第九、第十二章的关键材料。本章只回答一个问题：
\keyw{一只子供向的龙，如何成为一家有产业规模的 IP 公司}。

\section{起点：一家深圳公司}\label{sec:03:origin}

据《人民日报》2026 年 1 月 15 日报道，
2017 年，谢承恩与团队扎根深圳，
成立第七印象文化传媒（深圳）有限公司
\conf{A}\cite{jiang2026nailong}。
公司在 2019 年开始构思奶龙形象，其对外表述的创作定位是：
“想传递治愈、欢乐的积极情感，帮人们缓解焦虑、紧张等情绪”
（谢承恩语）\conf{A}\cite{jiang2026nailong}。
这一表述值得保留：它说明正典奶龙的\keyw{情感定位}
从构思期就锚定在“减压”而非“剧情深度”上——
后文将看到，这与奶蛙的“捧腹大笑”在功能上形成
一种耐人寻味的呼应（第六章、第七章详论）。

形象设计本身经历过迭代：团队“经过多次尝试，
最终设计出拥有黄色皮肤、呆萌可爱的奶龙形象”
\conf{A}\cite{jiang2026nailong}。
注意“黄色皮肤”这一细节：奶蛙保留了黄色，
而放弃了呆萌的圆脸比例——变种对正典的继承与偏离，
在配色层面就是“保留”，在比例层面就是“畸变”
（第四章考据、第六章符号分析的基础材料）。

\section{上线：短视频时代的冷启动}\label{sec:03:launch}

2020 年初，第一条奶龙短视频上线并迅速走红
\conf{A}\cite{jiang2026nailong}。
本书把其上线日定为 \textbf{2020 年 1 月 22 日}：
该日期来自六周年节点的倒推——
中国日报网 2026 年 1 月 28 日报道明确写道
“国产原创动漫 IP 奶龙，在 2026 年 1 月 22 日迎来了他的六周年庆”
\conf{B}\cite{chinadaily2026six}，
2026 减 6 年，恰为 2020-01-22。
本书将其标注为 \conf{B} 级（机构公开信息的日期倒推），
而非 \conf{A} 级，因为报道原文没有直接写出
“2020-01-22 上线”这一句。

上线后的早期扩张数据，目前只有第一版教材给出的口径：
“同年抖音粉丝破千万，单条视频 3 天点播破亿”
\conf{C}；2021 年，《奶龙》短视频合集在优酷、爱奇艺、腾讯
同步更新，有声读物专辑上线，“8 个月播放量破亿”\conf{C}。
这些数字无法在本书引用的 A/B 级来源中独立印证，
故全部保留 \conf{C} 标注，供修订者复核。

\begin{nlhistory}[史料与考据]
  2020 年短视频阶段的奶龙，与 2026 年的奶龙
  在商业形态上是两家公司：前者靠内容自然流量，
  后者靠“内容 + 授权”双轮驱动。谢承恩在人民日报报道中
  自述了这次转型：“形象要优化，内容要升级，
  商业化也还没起步”（2020 年初的瓶颈）；
  之后“团队从广告合作起步，逐步拓展到商业授权，
  积极探索‘内容 + 授权’双轮驱动的模式”\conf{A}
  \cite{jiang2026nailong}。
\end{nlhistory}

\section{剧集化：从短视频到番剧}\label{sec:03:series}

正典奶龙的剧集序列，各平台条目的编号方式并不完全一致
（有的平台把 2026 年的新番编作“第二季”的续作，
有的按系列名单列），本书按“系列名 + 上线时间”记录，
不强行统一季数编号：

\begin{table}[htbp]\centering
  \caption{正典奶龙剧集序列（以官方上线宣传与媒体公开信息为准）}
  \label{tab:03:series}
  \small
  \begin{tabularx}{\textwidth}{@{}llXc@{}}
    \toprule
    系列名 & 上线时间 & 备注 & 可信度 \\ \midrule
    《奶龙搞怪大作战》（第一季） & 2023 年
      & 进入高中、大学生视野的节点（第一版教材口径） & \conf{B/C} \\
    《奶龙与小七之大战暴暴龙》（第二季） & 2024 年
      & 在高中生与大学生之间“爆火”（第一版教材口径，未获独立权威源印证） & \conf{C} \\
    《奶龙之奇妙冒险》（第三季） & 2025 年 1 月
      & 一个月内全网播放量破 1 亿 & \conf{B} \\
    《奶龙搞怪大作战 2》 & 2026 年 1 月 28 日
      & 腾讯、优酷、爱奇艺等主流平台；覆盖风行、聚好看、
        银河少儿等大屏终端；嘉佳卡通、金鹰卡通等电视台陆续播出 & \conf{B} \\ \bottomrule
  \end{tabularx}
\end{table}

表中可见一条清晰的\keyw{渠道下沉}轨迹：
短视频平台（2020）→ 长视频平台同步（2021）→
番剧化（2023—）→ 大屏终端 + 电视台（2026）。
对一只子供向 IP 而言，大屏与电视台意味着
它进入了“家庭共同观看”的场域；而 2026 年 1 月 24 日
深圳海岸影城的六周年包场观影（约 200 位粉丝到场，
影片结束后“全场更自发齐唱生日歌”）
则说明其粉丝组织已具备线下动员能力
\conf{B}\cite{chinadaily2026six}。

\section{六周年：作为仪式的产业节点}\label{sec:03:anniversary}

2026 年 1 月 19 日至 25 日，第七印象举办“奶龙六周年”生日会，
含三大板块：线下观影会（新番抢先看）、五城地标大屏共庆
（武汉、青岛、天津、成都、苏州）、三城线下门店见面会
（武汉、乌鲁木齐、哈尔滨，联动茶饮品牌“吾饮良品”）
\conf{B}\cite{chinadaily2026six}。
1 月 22 日生日当天，奶龙以“歌手”身份发布首支单曲
《奶龙的小确幸》，生日直播掉落六周年限定周边与合作品牌礼包
\conf{B}\cite{chinadaily2026six}。
1 月 25 日，奶龙亮相中央广播电视总台（CMG）《动漫盛典》，
“以打头阵的高姿态惊艳亮相”（报道语）\conf{B}\cite{chinadaily2026six}。

\begin{figure}[htbp]\centering
\begin{tikzpicture}[font=\small]
  \draw[NLDeep,line width=1.2pt,-{Stealth[length=3mm]}] (-0.3,0) -- (14.1,0);
  \foreach \x in {0,1.7,3.4,6.1,7.8,9.5,11.6,13.2}{
    \fill[NLYellow] (\x,0) circle (2.2pt);
    \draw[NLDeep] (\x,0) -- (\x,0.12);
  }
  \node[font=\footnotesize,align=center,inner sep=1pt] at (0,0.95) {2017\\第七印象成立};
  \node[font=\footnotesize,align=center,inner sep=1pt] at (1.7,-0.95) {2019\\构思奶龙形象};
  \node[font=\footnotesize,align=center,inner sep=1pt] at (3.4,0.95) {2020-01-22\\短视频平台上线};
  \node[font=\footnotesize,align=center,inner sep=1pt] at (6.1,-0.95) {2023\\第一季《搞怪大作战》};
  \node[font=\footnotesize,align=center,inner sep=1pt] at (7.8,0.95) {2024\\第二季《大战暴暴龙》};
  \node[font=\footnotesize,align=center,inner sep=1pt] at (9.5,-0.95) {2025-01\\第三季《奇妙冒险》};
  \node[font=\footnotesize,align=left,inner sep=1pt] at (11.6,0.95) {2026-01-19/25\\六周年生日会\\（五城大屏 + 三城门店）};
  \node[font=\footnotesize,align=left,inner sep=1pt] at (13.2,-0.95) {2026-01-28\\《搞怪大作战 2》多平台上线};
\end{tikzpicture}
  \caption{正典奶龙产业里程碑（2017—2026）。2024 年节点可信度为
          \conf{C}（第一版教材口径），其余节点见事实手册。}
  \label{fig:03:timeline}
\end{figure}

这一周的意义，产业分析上叫“节点营销”，
文化分析上叫\keyw{年度仪式}：六周年庆典把“粉丝陪伴叙事”
（小红书“留言小纸条”活动，“数万条留言中，陪伴与成长
成为粉丝们提到的共同关键词”\conf{B}\cite{chinadaily2026six}）
转化为可见的消费与传播事件。对本书而言，它的价值在于
提供了一个\keyw{时间锚点}：奶蛙现象在 2026 年上半年全面爆发之前，
正典奶龙的官方动作集中在 1 月的这一周——
两个时间线的相对位置（官方仪式在前，民间爆火在后），
是第九章讨论“官方与二创关系”的基本史实依据。

\section{产业量化口径：为什么并列而非合并}\label{sec:03:numbers}

截至 2026 年初，两个最主要的来源给出了
数值不同但口径相近的量化描述。本书的处理方式如下表：

\begin{table}[htbp]\centering
  \caption{正典奶龙产业指标的两个来源口径（并列呈现，不作合并）}
  \label{tab:03:numbers}
  \small
  \begin{tabularx}{\textwidth}{@{}p{3.0cm}p{4.6cm}p{4.6cm}X@{}}
    \toprule
    指标 & 《人民日报》2026-01-15（\conf{A}）
      & 中国日报网 2026-01-28（\conf{B}） & 本书处理 \\ \midrule
    粉丝规模 & 全网粉丝超 3800 万 & 全球粉丝规模突破 3900 万+
      & 并列；“全网”与“全球”口径是否同构不明，不取均值 \\
    流量 & 流量破千亿 & 未提及 & 仅 A 级来源给出，保留 \\
    授权合作品牌 & 超 200 家 & 超 200 家（截至 2025 年底）
      & 两源一致，可作 A/B 级陈述 \\
    衍生产品 & 数超 2000 个 & SKU 逾 2000 个
      & 两源一致 \\
    销售地域 & 57 个国家和地区 & 未提及
      & 仅 A 级来源给出 \\
    微信表情包 & 未提及 & 累计发送量突破 80 亿次（截至 2025 年底）
      & 仅 B 级来源给出 \\
    译制 & 未提及 & 英语、韩语、印尼语等 15 种语言
      & 仅 B 级来源给出 \\
    政府支持 & 2024 年参评深圳市新型业态企业，获 25 万元扶持奖励
      & 未提及 & 仅 A 级来源给出 \\ \bottomrule
  \end{tabularx}
\end{table}

为什么本书不取 3800 万与 3900 万的均值？
三个理由。第一，\keyw{口径不明}：“全网”（国内平台合计？）
与“全球”（含海外平台？）的统计范围不同，
数值差异可能主要来自范围而非时间差。
第二，\keyw{时间差}：两报道相隔 13 天，
粉丝数的自然增长本身就能解释差距。
第三，\keyw{教学目的}：本书要训练的正是
“面对多个来源时不悄悄合并”的引用纪律——
取均值是一个看似中立、实则制造了不存在之精度的动作。
读者在第十二章的可复现性清单里会看到这一纪律的操作化版本。

\section{子供向作品为何溢出到大学生}\label{sec:03:overflow}

第一章提出的谜题是“为什么聪明人爱看无厘头”。
在正典一侧，这个谜题有一个更精确的版本：
奶龙的目标受众是\keyw{低龄观众}（“嘉佳卡通、金鹰卡通”等
播出终端\conf{B}\cite{chinadaily2026six}，
“连接亲子温情与成人解压的最大公约数”是报道原话
\conf{B}\cite{chinadaily2026six}），
但第一版教材的观察与高校同好会的存在都表明，
其实际受众显著外溢到了大学生群体
\conf{C}。本书给出三层解释，
均标注为本书分析（M 级），供检验而非供引用：

\begin{enumerate}[label=\arabic*.]
  \item \textbf{内容内核层}：奶龙的内容组织是“情绪单元”
        而非“叙事弧”——单个短视频/单集即可闭合，
        不需要追更成本。这与大学生碎片化注意力
        的接收条件匹配；而“治愈、欢乐”的定位
        （\conf{A}\cite{jiang2026nailong}）恰好对准了
        第一章所述的压力情境。
  \item \textbf{平台结构层}：短视频平台按兴趣分发，
        不按年龄分发。一个 20 岁用户看到奶龙，
        与一个 8 岁用户看到奶龙，在算法层面没有区别；
        “子供向”是内容标签，不是分发过滤器。
  \item \textbf{民间再生产层}：真正完成“大学生化”的
        不是官方内容，而是奶蛙式的二创——
        它把正典形象翻译成了大学生社群的
        内部语（空耳、抽象、立牌）。
        这一层在第四章与第九章展开。
\end{enumerate}

三层的分工是：第一层解释“为什么不排斥”，
第二层解释“为什么会看到”，第三层解释“为什么会主动传播”。
缺任何一层，都只能部分解释受众外溢。

\begin{nlcase}[案例 3-1：大屏与“全家桶”]
  2026 年 1 月，奶龙在武汉、天津、成都的极地海洋公园剧场大屏、
  青岛“极地之光·海上摩天轮”灯光秀、苏州阳澄半岛乐园
  “我是大明星剧场”同步播放六周年影像\conf{B}
  \cite{chinadaily2026six}。
  这些场景的共同特征是\keyw{家庭共同在场}
  （带孩子逛乐园的成年人）。
  一只子供向 IP 选择“五城地标大屏”而非“五城地铁广告”，
  其传播意图本身就承认了
  “观看者是亲子组合”这一结构——
  这为“受众外溢先于官方意图”的判断提供了
  一个反向参照：官方把资源投给了它认为的受众。
\end{nlcase}

\section{版权治理：官方如何对待“奶蛙式”二创}\label{sec:03:copyright}

产业线不能与法律线分开。据《人民日报》报道，
奶龙 IP 大热后，“侵权盗版问题”成为常态，
谢承恩列举的形态包括：“有商场未经授权搭建相关游乐设施，
IP 形象严重失真，安全性也无法保证；有人截取视频片段，
甚至用 AI 工具批量生成低质内容”\conf{A}\cite{jiang2026nailong}。
企业的应对是“组建专业法务团队”，
“在深圳市知识产权保护中心、市场监管部门等指导下，
依托地方快速维权机制与专项执法行动，及时维护了 IP 形象”
\conf{A}\cite{jiang2026nailong}；
海外则因“维权链条长、成本高”，
由广东省版权部门“协助对接海外知识产权组织与法律资源”
\conf{A}\cite{jiang2026nailong}。

与“奶蛙”直接相关的官方行为，目前只有 \conf{C} 级材料：
2026 年 3 月，“奶龙官方呼吁停止丑化二创，
关闭视频评论区并批量拉黑发布奶蛙的账号”
\conf{C}\cite{shoutanji2026naiwa}；
另一二手评述亦提及“奶龙官方将所有在奶龙视频评论区
中发奶蛙的人一律拉黑”，并称“直至今日你依然可以
在奶龙官方视频的评论区看见无数奶蛙图片”
\conf{C}\cite{shoutanji2026naiwa}。

这两条材料放在一起，读出一个对本书至关重要的结构：
\keyw{官方的抑制是不完全的，而不完全的抑制
恰好构成了迷因生态的缝隙}。第九章将论证，
如果 2026 年 3 月的关评与拉黑是彻底的，
奶蛙现象在时间上会被截断；
而它实际上没有——这一“失败的控制”
本身成为后续传播的燃料（第七章“逆反认证”机制）。
必须同时指出另一面：本章引用的全部 A 级材料
都说明官方维权是有组织、有法律依据的，
本书不对任何具体的二创行为作合法性判断（第十二章）。

\section{本章小结}\label{sec:03:summary}

本章回答了“正典奶龙如何成为产业”：
2017 年公司化、2019 年形象构思（定位“治愈、欢乐”）、
2020-01-22 上线（B 级日期倒推）、2021 年多平台 + 有声读物、
2023—2026 剧集化与大屏化、2026 年 1 月六周年仪式周、
以及“内容 + 授权”双轮驱动的规模化
（200+ 品牌、2000+ SKU、57 国、80 亿次表情包、15 种语言）。
在方法上，本章演示了两件事：
来源口径的并列呈现（表 \ref{tab:03:numbers}），
以及把“受众外溢”拆成内容内核、平台结构、
民间再生产三层可检验的解释。
遗留问题：正典一侧到此已经闭环，
但“变种”一侧的材料远没有这么齐——
奶蛙的起源至今是一条未结案的证据链。
这正是下一章的主题。

\begin{nlexercises}
  \item 表 \ref{tab:03:series} 中第二季（2024）标注为 \conf{C}
        而第三季（2025）标注为 \conf{B}：
        请用事实手册说明这一差异的依据，
        并解释为什么“同一系列、不同年份”可以
        拥有不同的可信度。
  \item 为“受众外溢三层解释”（\ref{sec:03:overflow}）
        各设计一条\keyw{可观察、可计数}的指标
        （例如：某平台 20 岁以上用户占比、
        二创作者中高校生占比、大屏场景的亲子观众比例），
        并指出哪一层最难测量、为什么。
  \item 官方“关评 + 拉黑”（\conf{C} 级）与
        “法务团队 + 快速维权机制”（\conf{A} 级）
        是否矛盾？请用第九章的框架
        （官方控制 vs 迷因生态缝隙）作答，
        并说明你的答案里哪些依赖 C 级材料、
        哪些只依赖 A 级材料。
\end{nlexercises}

\nlkeywords 第七印象、正典奶龙、剧集序列、六周年、
产业量化口径、受众外溢、版权治理、关评拉黑
